\emph{Semantik} kanggo basa tataran kapisan diwènèhaké déning
\emph{struktur} kanggo basa iku. Struktur ngemot \emph{ranah}
lan nemtokaké anggota ranah kanggo saben simbol konstanta.
Simbol fungsi diinterpretasi nganggo fungsi, lan simbol predikat
nganggo relasi ing ranah. Fungsi saka himpunan variabèl menyang
ranah diarani \emph{panugasan variabèl}. Relasi \emph{panyukupan}
nggandhèngaké struktur, panugasan variabèl, lan formula.
$\Sat{M}{!A}[s]$ ditetepaké kanthi induksi tumrap struktur~$!A$.
Asil iki mung gumantung marang interpretasi simbol kang pancèn
muncul ing~$!A$. Yèn $!A$ ora ngemot variabèl bebas, asilé ora
gumantung marang~$s$. Mula, yèn $!A$ minangka ukara,
$\Sat{M}{!A}$ yèn $\Sat{M}{!A}[s]$ kanggo sembarang~$s$.

Relasi panyukupan dadi dhasar kabèh konsèp semantis. Ukara~$!A$
\emph{sah}, $\Sat{{}}{!A}$, yèn disukupi ing saben struktur.
$!A$ minangka \emph{konsèkuènsi} saka himpunan ukara~$\Gamma$,
$\Gamma\Entails !A$, yèn lan mung yèn $\Sat{M}{!A}$ kanggo
saben $\Struct{M}$ kang nyukupi kabèh ukara ing~$\Gamma$.
$\Gamma$ \emph{bisa disukupi} yèn ana struktur kang nyukupi
saben ukara ing~$\Gamma$. Yèn ora, $\Gamma$ \emph{ora bisa disukupi}.
Konsèp-konsèp iki gegayutan. Upamané, $\Gamma\Entails !A$
yèn lan mung yèn $\Gamma\cup\{\lnot !A\}$ ora bisa disukupi.
